\documentclass[11pt,a4paper]{article}

\usepackage{fontspec}
\usepackage{polyglossia}
\setmainlanguage{russian}
\setmainfont{Liberation Serif}
\setsansfont{Liberation Sans}
\setmonofont{DejaVu Sans Mono}

\usepackage{amsmath,amssymb}
\usepackage{booktabs}
\usepackage{array}
\usepackage{geometry}
\geometry{margin=1.8cm}

\setlength{\parskip}{0.45em}
\setlength{\parindent}{0pt}

\begin{document}

\begin{center}
{\Large Ручной расчёт одной строки}\\[0.4em]
{\large Лабораторная работа №4}\\[0.2em]
Метод Симпсона с удвоением числа шагов
\end{center}

Сверяем программу \texttt{3\_n.py} с прямым счётом для одной точки таблицы.
Берём строку \(j = 1\). Это первая строка, где четырёх шагов уже не хватает
и программа переходит к восьми. На ней видно и саму формулу Симпсона, и правило остановки.

\section*{Что считаем}

Параметры из программы:
\[
a = 0,\quad b = 2,\quad c = 0,\quad d = 1,\quad m = 20,\quad \varepsilon = 0{,}001,
\]
\[
f(x, t) = \exp\left(-\left(\frac{t}{1 + x^{2}} + 0{,}001\, x\right)\right).
\]

Шаг по параметру \(t\):
\[
\tau = \frac{d - c}{m} = \frac{1}{20} = 0{,}05,
\qquad
t_{1} = c + 1\cdot\tau = 0{,}05.
\]

Нужна площадь
\[
F(0{,}05) = \int_{0}^{2} f(x,\, 0{,}05)\, dx.
\]

Формула Симпсона. Отрезок \([0, 2]\) делим на \(N\) равных частей длины \(h = 2/N\).
Число \(N\) обязательно чётное: программа начинает с \(N = 2\) и каждый раз удваивает его.
Узлы \(x_{i} = a + ih\), \(i = 0, 1, \ldots, N\). Крайние значения входят с коэффициентом \(1\),
нечётные внутренние узлы с коэффициентом \(4\), чётные внутренние с коэффициентом \(2\):
\[
S_{N} = \frac{h}{3}\Bigl(
f(x_{0}) + f(x_{N})
+ 4\sum_{k} f(x_{2k+1})
+ 2\sum_{k} f(x_{2k})
\Bigr).
\]
Внутренние суммы идут по тем \(i\), которые реально есть между \(0\) и \(N\).

Считаем \(S_{2}\), затем \(S_{4}\), затем \(S_{8}\).
Останавливаемся, когда соседние оценки впервые отличаются меньше чем на \(0{,}001\).
В ответ берём более мелкую из этих двух оценок. Так же делает программа.

Значения экспоненты округлены до 10 знаков после запятой.

\section*{Шаг 1. Два шага, \(N = 2\)}

\[
h = \frac{2}{2} = 1.
\]

Узлы: \(0,\ 1,\ 2\). Коэффициенты: \(1,\ 4,\ 1\).

\paragraph{Точка \(x = 0\).}
\[
\frac{0{,}05}{1 + 0} + 0 = 0{,}05,
\qquad
f(0) = e^{-0{,}05} = 0{,}9512294245.
\]
Коэффициент \(1\), вклад \(0{,}9512294245\).

\paragraph{Точка \(x = 1\).}
\[
1 + x^{2} = 2,
\qquad
\frac{0{,}05}{2} + 0{,}001\cdot 1 = 0{,}025 + 0{,}001 = 0{,}026,
\]
\[
f(1) = e^{-0{,}026} = 0{,}9743350896.
\]
Коэффициент \(4\), вклад \(4\cdot 0{,}9743350896 = 3{,}8973403584\).

\paragraph{Точка \(x = 2\).}
\[
1 + x^{2} = 5,
\qquad
\frac{0{,}05}{5} + 0{,}001\cdot 2 = 0{,}010 + 0{,}002 = 0{,}012,
\]
\[
f(2) = e^{-0{,}012} = 0{,}9880717129.
\]
Коэффициент \(1\), вклад \(0{,}9880717129\).

Сумма вкладов:
\[
0{,}9512294245 + 3{,}8973403584 + 0{,}9880717129 = 5{,}8366414958.
\]
Площадь:
\[
S_{2} = \frac{1}{3}\cdot 5{,}8366414958 = 1{,}9455471653.
\]

Сравнивать пока не с чем: это первая сетка. Запоминаем \(S_{2}\) и удваиваем число шагов.

\section*{Шаг 2. Четыре шага, \(N = 4\)}

\[
h = \frac{2}{4} = 0{,}5.
\]

Узлы: \(0;\ 0{,}5;\ 1;\ 1{,}5;\ 2\). Коэффициенты: \(1,\ 4,\ 2,\ 4,\ 1\).

{\small
\begin{center}
\begin{tabular}{@{}r r r r r@{}}
\toprule
\(x\) & аргумент & \(f(x)\) & коэф. & вклад \\
\midrule
0    & 0,05       & 0,9512294245 & 1 & 0,9512294245 \\
0,5  & 0,0405     & 0,9603091645 & 4 & 3,8412366580 \\
1    & 0,026      & 0,9743350896 & 2 & 1,9486701792 \\
1,5  & 0,0168846154 & 0,9832571308 & 4 & 3,9330285232 \\
2    & 0,012      & 0,9880717129 & 1 & 0,9880717129 \\
\bottomrule
\end{tabular}
\end{center}
}

Проверка одной строки, \(x = 0{,}5\):
\[
x^{2} = 0{,}25,\quad 1 + x^{2} = 1{,}25,\quad \frac{0{,}05}{1{,}25} = 0{,}04,
\]
\[
0{,}04 + 0{,}001\cdot 0{,}5 = 0{,}0405,
\qquad
e^{-0{,}0405} = 0{,}9603091645,
\]
\[
4\cdot 0{,}9603091645 = 3{,}8412366580.
\]

Сумма вкладов:
\[
0{,}9512294245 + 3{,}8412366580 + 1{,}9486701792 + 3{,}9330285232 + 0{,}9880717129
= 11{,}6622364978.
\]
Площадь:
\[
S_{4} = \frac{0{,}5}{3}\cdot 11{,}6622364978 = 1{,}9437060830.
\]

Сравнение с предыдущей сеткой:
\[
|S_{4} - S_{2}| = |1{,}9437060830 - 1{,}9455471653| = 0{,}0018410823.
\]
\(0{,}0018410823 > 0{,}001\). Заданная точность ещё не достигнута, число шагов снова удваиваем.

\section*{Шаг 3. Восемь шагов, \(N = 8\)}

\[
h = \frac{2}{8} = 0{,}25.
\]

Узлы от \(0\) до \(2\) с шагом \(0{,}25\). Коэффициенты по порядку:
\(1,\ 4,\ 2,\ 4,\ 2,\ 4,\ 2,\ 4,\ 1\).

{\small
\begin{center}
\begin{tabular}{@{}r r r r r@{}}
\toprule
\(x\) & аргумент & \(f(x)\) & коэф. & вклад \\
\midrule
0     & 0,05         & 0,9512294245 & 1 & 0,9512294245 \\
0,25  & 0,0473088235 & 0,9537927984 & 4 & 3,8151711936 \\
0,5   & 0,0405       & 0,9603091645 & 2 & 1,9206183290 \\
0,75  & 0,03275      & 0,9677804745 & 4 & 3,8711218980 \\
1     & 0,026        & 0,9743350896 & 2 & 1,9486701792 \\
1,25  & 0,0207621951 & 0,9794518553 & 4 & 3,9178074212 \\
1,5   & 0,0168846154 & 0,9832571308 & 2 & 1,9665142616 \\
1,75  & 0,0140576923 & 0,9860406557 & 4 & 3,9441626228 \\
2     & 0,012        & 0,9880717129 & 1 & 0,9880717129 \\
\bottomrule
\end{tabular}
\end{center}
}

Проверка одной строки, \(x = 0{,}25\):
\[
x^{2} = 0{,}0625,\quad 1 + x^{2} = 1{,}0625,
\]
\[
\frac{0{,}05}{1{,}0625} = 0{,}0470588235,
\qquad
0{,}001\cdot 0{,}25 = 0{,}00025,
\]
\[
0{,}0470588235 + 0{,}00025 = 0{,}0473088235,
\]
\[
e^{-0{,}0473088235} = 0{,}9537927984,
\qquad
4\cdot 0{,}9537927984 = 3{,}8151711936.
\]

Сумма всех девяти вкладов равна \(23{,}3233670428\). Площадь:
\[
S_{8} = \frac{0{,}25}{3}\cdot 23{,}3233670428 = 1{,}9436139202.
\]

Сравнение с сеткой из четырёх шагов:
\[
|S_{8} - S_{4}| = |1{,}9436139202 - 1{,}9437060830| = 0{,}0000921628.
\]
\(0{,}0000921628 < 0{,}001\). Критерий выполнен, дальше удваивать не нужно.

\section*{Ответ и сверка с программой}

\[
F(0{,}05) \approx S_{8} = 1{,}9436139202, \qquad N = 8.
\]

В таблицу пишется \(S_{8}\), а не \(S_{4}\): берётся сумма на той сетке, где условие уже выполнилось.
Программа печатает значение с десятью знаками после запятой:

\begin{center}
\texttt{1 | 0.0500 | 1.9436139202 | 8}
\end{center}

Ручной счёт даёт те же \(t_{1} = 0{,}05\), \(F = 1{,}9436139202\) и \(N = 8\).

Почему именно 8, а не 4: разность \(S_{4}\) и \(S_{2}\) ещё больше допуска \(0{,}001\)
(получилось \(0{,}00184\)). Разность \(S_{8}\) и \(S_{4}\) уже меньше допуска
(получилось \(0{,}000092\)).

\end{document}
